\textit{Gleichung, wo die gesuchte Grösse eine Funktion ist. Ableitung der gesuchten Funktion kann vorkommen.}

\subsection{Aufstellen von Differentialgleichungen}

\textit{\textbf{Typische Situation:} Gegeben sind verschiedenen Informationen und gesucht ist die
Entwicklung einer bestimmten Grösse.}

\begin{enumerate}
    \item Wie verändert sich die betrachtete Grösse in einem kleinen Zeitintervall $\delta t$?\\
        \textit{Zuwachs / Abnahme / Nettobilanz}
    \item Was ist die Differenz zwischen den Werten, falls die Zeiten nahe beieinander sind?
    \item Was ist die proportionale Veränderung? \textit{(Quotient)}
    \item Was passiert mit dem Grenzwert $\delta t \rightarrow 0$?
    \item Zusammenfassen 
\end{enumerate}

\textit{
    Für jede Kraft $F$ gilt insbesondere (nach Newton) $F = m\ddot{y} = ma$
}

\subsection{Definition / Klassifizierung}

Eine \definition{Differentialgleichung} ist eine Gleichung, in der die gesuchte Funktion, sowie deren Ableitungen auftreten.

\textit{\textbf{Beispiele:} $u'(t)=f(a-u(t)), u'(t)=r(a-u(t))(b-u(t)), u'(t)=ru(t)(1-\frac{u(t)}{L})$ wobei $r,a,b \text{ und }L$ Konstanten sind, oder auch $y^{(4)}(x)-4y^{(3)}(x)+5y^{(2)}(x)-4y'(x)+4y(x)=0$}

\textit{
$t,x$ heissen \textbf{unabhängige Variablen},
$u,y$ heissen \textbf{abhängige Variablen}
}

Eine Differentialgleichung der Form $a_n(x)y^{(n)}(x)+a_{n-1}(x)y^{(n-1)}(x)+\dots+a_1(x)y'(x)+a_0(x)y(x)=s(x)$ heisst \definition{lineare Differentialgleichung} mit \textbf{Koeffizienten $a_i(t)$}.

Die \definition{Ordnung} der Differentialgleichung ist die maximale Ordnung der vorkommenden Ableitungen.

Die Funktion $s(x)$ nennen wir \definition{Störfunktion.}

Falls gilt $s(x)=0$ heisst die Differentialgleichung \definition{homogen}, andernfalls \definition{inhomogen}.

Falls alle Funktionen $a_i(x)$ in der Differentialgleichung konstante Funktionen sind, nennen wir diese eine \definition{Differentialgleichung mit konstanten Koeffizienten}.

\theorem{Superpositionsprinzip}
Sind $y_1(x)$ und $y_2(x)$ Lösungen einer linearen, homogenen Differentialgleichung, so ist auch $y(x)=C_1y_1(x)+C_2y_2(x), C_1, C_2 \in \mathbb{R}$ eine Lösung derselben Differentialgleichung.

\textit{Geht nur für homogene Differentialgleichungen (weil falls $C_1=C_2=0$)}

\theorem{Fundamentallösungen}
Eine lineare, homogene Differentialgleichung der Ordnung $n$ besitzt $n$ Lösungen $y_1(x), \dots, y_n(x)$, so ist die allgemeine Lösung dieser Gleichung gegeben durch $y(x)=C_1y_1(x)+\dots+C_ny_n(x), C_1, \dots, C_n \in \mathbb{R}$.
Diese Lösungen werden \definition{Fundamentallösungen} genannt.

\begin{exercise}
    Zeige dass alle Funktionen der Form $y(x)=a \cosh (\frac{x-x_0}{a})+y_0$ die Differentailgleichung $(y-y_0)y''-(y')^2 = 1$ lösen:

    \textit{$y(x)$ "berechnen": Ableitungen berechnen, einsetzen und rechnerisch nachweisen, dass dies $=1$.}

    \textit{Für Anfangswertproblem die Anfangsbedingungen ebenfalls überprüfen.}
\end{exercise}

\subsection{Lösen einer linearen, homogenen Differentialgleichung}
$$a_n u^{(n)}(x) + a_{n-1} u^{(n-1)}(x) \dots + a_1 u'(x) + a_0 u(x) = 0$$

\definition{Eulerscher Ansatz} $u(x)=e^{\lambda x}$

\definition{Charakteristische Gleichung} $a_n \lambda^n + a_{n-1} \lambda^{n-1}+ \dots + a_1 \lambda + a_0 =0$
\textit{(folgt aus dem eulerschen Ansatz)} 

\definition{Charakteristisches Polynom} $p(\lambda) = a_n \lambda^n + a_{n-1} \lambda^{n-1}+ \dots + a_1 \lambda + a_0$.

\textit{
    Jede einfache relle Nullstelle $\lambda$ des charakteristischen Polynoms erzeugt die Fundamentallösung $u(x)=e^{\lambda x}$\\
    Jede Paar zweier zueinande komplex konjugierter Nullstellen $\lambda_{1,2}=a\pm ib$ erzeugt die beiden Fundamentallösungen $u_1(x)=e^{ax}\sin(bx)$, $u_2(x)=e^{ax}\cos(bx)$\\
    Jede Nullstelle $\lambda$ mit Multiplizität $s$ erzeugt die $s$ Fundamentallösungen $x^ke^{\lambda x}, k=0,...,s-1$
}

\greenBox{Allgemeine Lösung:}
\begin{itemize}
    \item Falls alle $r$ Nullstellen $\lambda_i$ des charakteristischen Polynoms reell sind und Multiplizität $m_i$ haben, gilt $u(x)=\sum_{i=0}^{r}(\sum_{p=0}^{m_i - 1}C_{ip} x^p e^{\lambda_i x})$
    \item Falls wir $r$ reelle Nullstellen und $s$ \textbf{Paare} komplexer Nullstellen $(a_j \pm ib_j)$ mit jeweils Multiplizität $m_i$ haben, gilt $u(x)=\sum_{i=0}^{r}(\sum_{p=0}^{m_i - 1}C_{ip} x^p e^{\lambda_i x}) + \sum_{j=0}^{s} (\sum_{q=0}^{m_j - 1}(A_{jq} x^q e^{a_j x} \sin(b_j x)+B_{jq} x^q e^{a_j x} \cos(b_j x)))$    
\end{itemize}

\begin{exercise}
    Allgemeine Lösung einer linearen, homogenen Differentialgleichung mit Parametern, bspw. $u''(x)-(a+b)u'(x)+abu(x)=0$

    \textit{Charakteristisches Polynom lösen; ergibt $\lambda_1 = a, \lambda_2 = b$, dann Fallunterscheidung für $a\neq b$ und $a=b$. Separat allgemeine Lösungen bestimmen (siehe oben).}

    \small{\textit{Wenn man eine Anfangswertaufgabe mit dieser Gleichung anschaut, kann man mithilfe von Taylorreihen und Vereinfachen zeigen, dass die Lösung für $a \neq b$ und $b \rightarrow a$ gegen die Lösung für den Fall $a=b$ konvergiert.}}
\end{exercise}

\subsection{Lösen von Differentialgleichungen erster Ordnung}
\begin{itemize}
    \item Trennung der Variablen
    \item Substitution
    \item Variation der Konstanten (lineare, inhomogene Differentialgleichung)
\end{itemize}

\subsection{Lösen von linearen, inhomogenen Differentialgleichung}
\begin{itemize}
    \item Geeignete Ansätze für partikuläre Lösung
\end{itemize}

\subsection{Lösen anderer Differentialgleichungen}

% TODO